%! Unit 4: Rational Functions — Summative Assessment — Practice Test --- Study Copy (blank)
\documentclass[10pt]{article}
\usepackage{algebra2-article}
\usepackage{algebra2-boxes}
\usepackage{grading}
\gradingassignment{alg2-u04-practice-test}
\gradingcourse{algebra_2}
\usepackage{algebra2-graded}
% Coordinate grid: {xmin}{xmax}{ymin}{ymax}
\newcommand{\lgrid}[4]{%
  \draw[step=1, linegray!40, very thin] (#1,#3) grid (#2,#4);
  \draw[->, semithick, charcoal] (#1,0)--(#2,0) node[below right, font=\scriptsize]{$x$};
  \draw[->, semithick, charcoal] (0,#3)--(0,#4) node[left, font=\scriptsize]{$y$};}
\newcommand{\dpt}[2]{\fill[charcoal] (#1,#2) circle (3pt);}

% Part-divider strip (headlinebox takes one arg: the background color)
\newcommand{\parthead}[1]{%
  \vspace{6pt}\noindent
  \begin{headlinebox}{forest}\color{white}\bfseries #1\end{headlinebox}%
  \vspace{4pt}}

\begin{document}

\pageheader{Unit 4: Rational Functions}{Practice Test --- Study Copy}

\vspace{0.10in}
\begin{remindbox}
\small
\textbf{This is a practice test.} It mirrors the real Unit 4 test in length, format,
and the ideas it covers, but the numbers and contexts differ. Work every problem
with no notes, then check your answers against the key.
\end{remindbox}

% ======================================================================
\parthead{Part A --- Vocabulary (8 pts)}
Write the letter of the best description next to each term.

\vspace{4pt}
\noindent
\begin{minipage}[t]{0.40\textwidth}
\begin{enumerate}[leftmargin=2.2em, itemsep=7pt, label=\arabic*.]
  \item Rational function \inblank{\blank{1.2cm}}{\gradedblank[id=qa1, type=short, points=1, answer={D}, check=text]{1.2cm-2pt}{10pt}}
  \item Domain restriction \inblank{\blank{1.2cm}}{\gradedblank[id=qa2, type=short, points=1, answer={A}, check=text]{1.2cm-2pt}{10pt}}
  \item Vertical asymptote \inblank{\blank{1.2cm}}{\gradedblank[id=qa3, type=short, points=1, answer={H}, check=text]{1.2cm-2pt}{10pt}}
  \item Hole \inblank{\blank{1.2cm}}{\gradedblank[id=qa4, type=short, points=1, answer={E}, check=text]{1.2cm-2pt}{10pt}}
  \item Horizontal asymptote \inblank{\blank{1.2cm}}{\gradedblank[id=qa5, type=short, points=1, answer={B}, check=text]{1.2cm-2pt}{10pt}}
  \item Least common denominator \inblank{\blank{1.2cm}}{\gradedblank[id=qa6, type=short, points=1, answer={G}, check=text]{1.2cm-2pt}{10pt}}
  \item Extraneous solution \inblank{\blank{1.2cm}}{\gradedblank[id=qa7, type=short, points=1, answer={F}, check=text]{1.2cm-2pt}{10pt}}
  \item Constant of variation \inblank{\blank{1.2cm}}{\gradedblank[id=qa8, type=short, points=1, answer={C}, check=text]{1.2cm-2pt}{10pt}}
\end{enumerate}
\end{minipage}%
\hfill
\begin{minipage}[t]{0.57\textwidth}
\small
\begin{description}[leftmargin=1.4em, itemsep=3pt]
  \item[(A)] A value that makes the \emph{original} denominator zero, so it can never be an input.
  \item[(B)] A horizontal line the graph approaches at its far left and far right; found by comparing the degrees of the numerator and denominator.
  \item[(C)] The number $k$ in $y=kx$ or $y=\frac{k}{x}$; you find it from one known pair of values.
  \item[(D)] A quotient of two polynomials whose denominator is not zero.
  \item[(E)] A single missing point on an otherwise unbroken curve, left behind by a factor that cancels.
  \item[(F)] A number that turns up when you solve but must be thrown out, because it is an excluded value of the original equation.
  \item[(G)] The smallest expression that every denominator divides evenly, built from the factored denominators.
  \item[(H)] A vertical line the graph runs alongside but never touches, left by a factor that does \emph{not} cancel.
\end{description}
\end{minipage}

% ======================================================================
\parthead{Part B --- Multiple Choice (12 pts)}
Circle the best answer.

\begin{enumerate}[leftmargin=1.9em, itemsep=9pt]
  \item The domain of $f(x)=\dfrac{x+5}{x^2-9}$ is all real numbers except
    \hfill (A) $x=-5$ \quad (B) $x=3$ \quad (C) $x=\pm 3$ \quad (D) none --- all reals\nobreak\hfill\llap{\raisebox{0pt}[0pt][0pt]{\small\textbf{Letter:}\ \gradedblank[id=qb1, type=choice, points=2, answer={C}]{0.4in}{10pt}}}
  \item Simplified completely, $\dfrac{x^2-25}{x^2+3x-10}$ is
    \hfill (A) $\dfrac{x-5}{x-2}$, $x\ne 2$ \quad (B) $\dfrac{x-5}{x-2}$, $x\ne -5,\,2$ \quad (C) $\dfrac{x+5}{x-2}$, $x\ne 2$ \quad (D) $\dfrac{-5}{3x-10}$\nobreak\hfill\llap{\raisebox{0pt}[0pt][0pt]{\small\textbf{Letter:}\ \gradedblank[id=qb2, type=choice, points=2, answer={B}]{0.4in}{10pt}}}
  \item The horizontal asymptote of $f(x)=\dfrac{3x^2-1}{x^2+4x}$ is
    \hfill \llap{\hbox{\gradedblank[id=qb3, type=choice, points=2, answer={C}]{0.2in}{10pt}}\,}(A) $y=0$ \quad (B) $y=\frac13$ \quad (C) $y=3$ \quad (D) there is none
  \item \letterbox{\gradedblank[id=qb4, type=choice, points=2, answer={B}]{0.4in}{10pt}}Which statement about $h(x)=\dfrac{x^2-x-6}{x^2-9}$ is true?
    \\[3pt]\hspace*{1.2em} (A) There are vertical asymptotes at $x=3$ and at $x=-3$.
    \\[2pt]\hspace*{1.2em} (B) There is a hole at $x=3$ and a vertical asymptote at $x=-3$.
    \\[2pt]\hspace*{1.2em} (C) There is a hole at $x=-3$ and a vertical asymptote at $x=3$.
    \\[2pt]\hspace*{1.2em} (D) There is a hole at $x=-2$ and a vertical asymptote at $x=3$.
  \item The graph of $g(x)=\dfrac{1}{x-4}+2$ has asymptotes
    \hfill (A) $x=4$, $y=2$ \quad (B) $x=-4$, $y=2$ \quad (C) $x=4$, $y=0$ \quad (D) $x=2$, $y=4$\nobreak\hfill\llap{\raisebox{0pt}[0pt][0pt]{\small\textbf{Letter:}\ \gradedblank[id=qb5, type=choice, points=2, answer={A}]{0.4in}{10pt}}}
  \item $y$ varies \emph{inversely} as $x$, and $y=8$ when $x=3$. When $x=12$, $y=$
    \hfill \llap{\hbox{\gradedblank[id=qb6, type=choice, points=2, answer={A}]{0.35in}{10pt}}\enspace}(A) $2$ \quad (B) $6$ \quad (C) $24$ \quad (D) $32$
\end{enumerate}

% ======================================================================
\parthead{Part C --- Short Answer \& Computation (40 pts)}
Show your work. State every restriction on the variable.

\begin{enumerate}[leftmargin=1.9em, itemsep=10pt]
  \item \textbf{Read the graph.} The rational function $f$ is graphed below; the dashed lines are its
        asymptotes and each gridline is one unit.
        \begin{center}
        \begin{tikzpicture}[scale=0.45]
          \lgrid{-6}{6}{-6}{6}
          \draw[navy, dashed, semithick] (2,-6)--(2,6);
          \draw[navy, dashed, semithick] (-6,1)--(6,1);
          \begin{scope}
            \clip (-6,-6) rectangle (6,6);
            \draw[very thick, forest, domain=-6:1.14, samples=160, smooth]
              plot (\x,{((\x)+4)/((\x)-2)});
            \draw[very thick, forest, domain=3.2:6, samples=160, smooth]
              plot (\x,{((\x)+4)/((\x)-2)});
          \end{scope}
          \foreach \x in {-4,-2,2,4}
            \draw[charcoal] (\x,-0.15)--(\x,0.15) node[below=1pt, font=\tiny]{$\x$};
          \dpt{-4}{0}
          \dpt{0}{-2}
        \end{tikzpicture}
        \end{center}
        (a) Equation of the vertical asymptote: \inblank{\blank{2.0cm}}{\gradedblank[id=qc1a, type=short, points=1, answer={x=2}]{2.0cm-2pt}{10pt}} \quad
            Equation of the horizontal asymptote: \inblank{\blank{2.0cm}}{\gradedblank[id=qc1b, type=short, points=1, answer={y=1}]{2.0cm-2pt}{10pt}}
        \\[4pt]
        (b) $x$-intercept: \inblank{\blank{2.0cm}}{\gradedblank[id=qc1c, type=short, points=1, answer={(-4,0)}, check=text]{2.0cm-2pt}{10pt}}; \quad $y$-intercept: \inblank{\blank{2.0cm}}{\gradedblank[id=qc1d, type=short, points=1, answer={(0,-2)}, check=text]{2.0cm-2pt}{10pt}}
        \\[4pt]
        (c) Domain (interval notation): \inblank{\blank{4.2cm}}{\gradedblank[id=qc1e, type=short, points=1, answer={(-inf,2)U(2,inf)}, check=text, display={$(-\infty,2)\cup(2,\infty)$}]{4.2cm-2pt}{10pt}}
        \\[4pt]
        (d) Range (interval notation): \inblank{\blank{4.2cm}}{\gradedblank[id=qc1f, type=short, points=1, answer={(-inf,1)U(1,inf)}, check=text, display={$(-\infty,1)\cup(1,\infty)$}]{4.2cm-2pt}{10pt}}
        \\[4pt]
        (e) End behavior: as $x\to\pm\infty$, $f(x)\to$ \inblank{\blank{1.6cm}}{\gradedblank[id=qc1g, type=short, points=1, answer={1}, check=numeric]{1.6cm-2pt}{10pt}}
        \\[4pt]
        (f) Interval(s) where $f$ is decreasing: \inblank{\blank{4.6cm}}{\gradedblank[id=qc1h, type=short, points=1, answer={(-inf,2), (2,inf)}, check=set, display={$(-\infty,2)$ and $(2,\infty)$ --- decreasing on \emph{each} branch, but not on their union}]{4.6cm-2pt}{10pt}}
        \\[4pt]
        (g) Interval(s) where $f(x)>0$: \inblank{\blank{4.6cm}}{\gradedblank[id=qc1i, type=short, points=1, answer={(-inf,-4)U(2,inf)}, check=text, display={$(-\infty,-4)\cup(2,\infty)$}]{4.6cm-2pt}{10pt}}

  \item \textbf{Simplify and state the restrictions.} Read the restrictions from the \emph{original}
        denominator.
        \\[4pt]
        (a) $\dfrac{3x^2-27}{x^2+x-12}$ \hspace{2.5cm} (b) $\dfrac{x^2-16}{4-x}$
        \vspace{2.8cm}\vadjust{\workanswer[(a)]{\gradedblank[id=qc2a, type=short, points=1, answer={3(x+3)/(x+4)}, display={$3(x+3)/(x+4)$}]{1.3in}{10pt}\ \textbf{$x\ne$}\ \gradedblank[id=qc2b, type=short, points=1, answer={3, -4}, check=set, display={$3,\ -4$}]{0.8in}{10pt}\quad\textbf{(b)}\ \gradedblank[id=qc2c, type=short, points=1, answer={-(x+4)}]{1.3in}{10pt}\ \textbf{$x\ne$}\ \gradedblank[id=qc2d, type=short, points=1, answer={4}, check=set]{0.8in}{10pt}}}

  \item \textbf{Multiply or divide.} Factor first; give the simplest form with its restrictions.
        \\[6pt]
        (a) $\dfrac{x^2-4}{x^2+6x+9}\cdot\dfrac{x+3}{x-2}$
        \vspace{2.4cm}\vadjust{\workanswer{\gradedblank[id=qc3a, type=short, points=1, answer={(x+2)/(x+3)}, display={$(x+2)/(x+3)$}]{1.3in}{10pt}\ \textbf{$x\ne$}\ \gradedblank[id=qc3b, type=short, points=1, answer={-3, 2}, check=set, display={$-3,\ 2$}]{0.8in}{10pt}}}
        \\
        (b) $\dfrac{x^2-x-6}{x^2-1}\div\dfrac{x-3}{x+1}$
        \vspace{2.6cm}\vadjust{\workanswer{\gradedblank[id=qc3c, type=short, points=1, answer={(x+2)/(x-1)}, display={$(x+2)/(x-1)$}]{1.3in}{10pt}\ \textbf{$x\ne$}\ \gradedblank[id=qc3d, type=short, points=1, answer={1, -1, 3}, check=set, display={$1,\ -1,\ 3$}]{0.9in}{10pt}}}

  \item \textbf{Add, subtract, and simplify a complex fraction.}
        \\[6pt]
        (a) $\dfrac{5}{x+2}-\dfrac{3}{x-1}$
        \vspace{2.2cm}\vadjust{\workanswer{\gradedblank[id=qc4a, type=short, points=1, answer={(2x-11)/((x+2)(x-1))}, display={$(2x-11)/\bigl((x+2)(x-1)\bigr)$}]{1.6in}{10pt}\ \textbf{$x\ne$}\ \gradedblank[id=qc4b, type=short, points=1, answer={-2, 1}, check=set, display={$-2,\ 1$}]{0.8in}{10pt}}}
        \\
        (b) $\dfrac{x}{x^2-9}+\dfrac{2}{x+3}$
        \vspace{2.4cm}\vadjust{\workanswer{\gradedblank[id=qc4c, type=short, points=1, answer={3(x-2)/((x-3)(x+3))}, display={$3(x-2)/\bigl((x-3)(x+3)\bigr)$}]{1.6in}{10pt}\ \textbf{$x\ne$}\ \gradedblank[id=qc4d, type=short, points=1, answer={3, -3}, check=set, display={$\pm 3$}]{0.8in}{10pt}}}
        \\
        (c) $\dfrac{\ \frac{1}{x}+\frac{1}{2}\ }{\ \frac{1}{x}-\frac{1}{2}\ }$
        \vspace{2.6cm}\vadjust{\workanswer{\gradedblank[id=qc4e, type=short, points=1, answer={(2+x)/(2-x)}, display={$(2+x)/(2-x)$}]{1.3in}{10pt}\ \textbf{$x\ne$}\ \gradedblank[id=qc4f, type=short, points=1, answer={0, 2}, check=set, display={$0,\ 2$}]{0.8in}{10pt}}}

  \item \textbf{The parent function and its transformations.} Let $g(x)=\dfrac{1}{x+3}-2$.
        \\[4pt]
        (a) Describe the transformation(s) that take $y=\dfrac1x$ to $g$. \inblank{\blank{6.0cm}}{\gradedblank[id=qc5a, type=short, points=1, answer={left 3, down 2}, check=text, display={left 3 units, down 2 units}]{6.0cm-2pt}{10pt}}
        \\[4pt]
        (b) Vertical asymptote: \inblank{\blank{1.8cm}}{\gradedblank[id=qc5b, type=short, points=1, answer={x=-3}]{1.8cm-2pt}{10pt}}; \quad horizontal asymptote: \inblank{\blank{1.8cm}}{\gradedblank[id=qc5c, type=short, points=1, answer={y=-2}]{1.8cm-2pt}{10pt}}
        \\[4pt]
        (c) Domain: \inblank{\blank{2.4cm}}{\gradedblank[id=qc5d, type=short, points=1, answer={x!=-3}, display={$x\ne -3$}]{2.4cm-2pt}{10pt}}; \quad range: \inblank{\blank{2.4cm}}{\gradedblank[id=qc5e, type=short, points=1, answer={y!=-2}, display={$y\ne -2$}]{2.4cm-2pt}{10pt}}
        \\[4pt]
        (d) A different hyperbola of the form $y=\dfrac{1}{x-h}+k$ has vertical asymptote $x=1$ and
            horizontal asymptote $y=4$. Write its equation. \inblank{\blank{3.6cm}}{\gradedblank[id=qc5f, type=short, points=1, answer={y=1/(x-1)+4}, display={$y=1/(x-1)+4$}]{3.6cm-2pt}{10pt}}

  \item \textbf{Analyze from the equation.} For $f(x)=\dfrac{2x^2-8}{x^2-x-6}$:
        \\[4pt]
        (a) Factor the numerator and denominator, and list every excluded value.
        \vspace{1.8cm}\vadjust{\workanswer[Excluded:]{\gradedblank[id=qc6a, type=short, points=1, answer={3, -2}, check=set, display={$3,\ -2$}]{1.0in}{10pt}\quad\textbf{(b) Hole:}\ \gradedblank[id=qc6b, type=short, points=1, answer={(-2,8/5)}, check=text, display={$(-2,8/5)$}]{1.0in}{10pt}}}
        \\
        (b) Which excluded value gives a \emph{hole} and which gives a \emph{vertical asymptote}?
            Give the coordinates of the hole.
        \vspace{2.0cm}
        \\
        (c) Horizontal asymptote: \inblank{\blank{2.0cm}}{\gradedblank[id=qc6c, type=short, points=1, answer={y=2}]{2.0cm-2pt}{10pt}} \quad (Justify with a degree comparison.)
        \\[4pt]
        (d) $x$-intercept: \inblank{\blank{2.0cm}}{\gradedblank[id=qc6d, type=short, points=1, answer={(2,0)}, check=text]{2.0cm-2pt}{10pt}}; \quad $y$-intercept: \inblank{\blank{2.0cm}}{\gradedblank[id=qc6e, type=short, points=1, answer={(0,4/3)}, check=text, display={$(0,4/3)$}]{2.0cm-2pt}{10pt}}

  \item \textbf{Solve. Check every answer against the excluded values.}
        \\[6pt]
        (a) $\dfrac{x}{x+2}+\dfrac{1}{x-1}=1$
        \vspace{2.6cm}\vadjust{\workanswer[Solution:]{\gradedblank[id=qc7a, type=short, points=1, answer={4}, check=set]{1.0in}{10pt}}}
        \\
        (b) $\dfrac{x}{x-4}=\dfrac{4}{x-4}+3$
        \vspace{2.6cm}\vadjust{\workanswer[Solution:]{\gradedblank[id=qc7b, type=short, points=1, answer={no solution}, check=text]{1.0in}{10pt}}}
        \\
        (c) $\dfrac{x}{x-2}+\dfrac{2}{x}=\dfrac{4}{x^2-2x}$
        \vspace{3.0cm}\vadjust{\workanswer[Solution:]{\gradedblank[id=qc7c, type=short, points=1, answer={-4}, check=set]{1.0in}{10pt}}}

  \item \textbf{Variation.} The table below records a pair of related quantities.
        \begin{center}
        \renewcommand{\arraystretch}{1.25}
        \begin{tabular}{|c|c|c|c|c|}
        \hline
        $x$ & $2$ & $4$ & $5$ & $10$ \\ \hline
        $y$ & $30$ & $15$ & $12$ & $6$ \\ \hline
        \end{tabular}
        \end{center}
        (a) Is $y$ directly or inversely proportional to $x$? Give the computation from the table that
            proves it.
 \rubricbox{\gradedblank[id=qc8a, type=rubric, points=1, rubric={../rubrics/alg2-u04-practice-test-qc8a.yaml}, display={Inversely --- the \emph{products} are constant: $2\cdot30=4\cdot15=5\cdot12=10\cdot6=60$ (the ratios $y/x$ are not)}]{\linewidth}{10pt}}
        \\[4pt]
        (b) Find the constant of variation $k$ and write the equation. \inblank{\blank{5.0cm}}{\gradedblank[id=qc8b, type=short, points=1, answer={y=60/x}, display={$k=60$, so $y=60/x$}]{5.0cm-2pt}{10pt}}
        \\[4pt]
        (c) Use your equation to find $y$ when $x=8$. \inblank{\blank{3.0cm}}{\gradedblank[id=qc8c, type=short, points=1, answer={7.5}, check=numeric, display={$y=60/8=7.5$}]{3.0cm-2pt}{10pt}}
\end{enumerate}

% ======================================================================
\parthead{Part D --- Extended Response (12 pts)}
Explain your reasoning in full sentences.

\begin{enumerate}[leftmargin=1.9em, itemsep=16pt]
  \item \textbf{The full analysis.} Consider $f(x)=\dfrac{x^2-5x+4}{x^2-16}$.
        \textbf{(a)} Factor both parts and list every excluded value of the \emph{original} function.
        \textbf{(b)} Run the cancel test: which excluded value is a hole and which is a vertical
        asymptote? Give the hole's coordinates and explain how you found its height.
        \textbf{(c)} State the horizontal asymptote, justify it with a degree comparison, and write the
        end behavior it describes.
        \textbf{(d)} Give the $x$- and $y$-intercepts.
        \textbf{(e)} Write the domain in interval notation.
        \textbf{(f)} A classmate simplifies $f$ to $\dfrac{x-1}{x+4}$ and says the two functions are the
        same. Say exactly where they agree and where they do not.
        \rubricbox{\gradedblank[id=qd1, type=rubric, points=6, rubric={../rubrics/alg2-u04-practice-test-qd1.yaml}, display={(a) $f(x)=\dfrac{(x-1)(x-4)}{(x-4)(x+4)}$. The original denominator is zero at $x=4$ and
        $x=-4$, so both are excluded.\\[3pt]
        (b) $(x-4)$ cancels, so $x=4$ is a \emph{hole}; $(x+4)$ does not, so $x=-4$ is a
        \emph{vertical asymptote}. The hole's height comes from the simplified form
        $\dfrac{x-1}{x+4}$ evaluated at $x=4$: $\dfrac{3}{8}$. The hole is
        $\left(4,\tfrac38\right)$. Substituting $4$ into the \emph{original} would give $\tfrac00$,
        which is why the simplified form is the one to use.\\[3pt]
        (c) Both degrees are $2$, so the horizontal asymptote is the ratio of leading coefficients,
        $\tfrac11=1$: the line $y=1$. That is the end-behavior statement --- as $x\to-\infty$ and as
        $x\to+\infty$, $f(x)\to 1$.\\[3pt]
        (d) $x$-intercept: the simplified numerator is zero at $x=1$, giving $(1,0)$. $y$-intercept:
        $f(0)=\dfrac{4}{-16}=-\dfrac14$, giving $\left(0,-\tfrac14\right)$.\\[3pt]
        (e) $(-\infty,-4)\cup(-4,4)\cup(4,\infty)$.\\[3pt]
        (f) They agree at every input both accept --- every $x$ other than $4$ and $-4$. They do not
        agree at $x=4$: the simplified expression happily returns $\tfrac38$ there, but $f(4)$ does
        not exist. Two equivalent expressions can have different \emph{domains}, so cancelling never
        restores an input.}]{\linewidth}{6.0cm-6pt}}

  \item \textbf{Modeling.} A club pays a one-time setup fee of \$180 to a screen printer, plus \$7 for
        each shirt printed. The \emph{average cost per shirt} for an order of $n$ shirts is
        \[ A(n)=\frac{7n+180}{n}. \]
        \textbf{(a)} Explain why $A$ is a rational function, and state the domain that makes sense for
        this situation.
        \textbf{(b)} Find $A(20)$ and $A(60)$, with units. What is happening to the average cost?
        \textbf{(c)} Rewrite $A(n)$ in the form $7+\dfrac{180}{n}$ and use it to state the horizontal
        asymptote. What does that number mean about this order?
        \textbf{(d)} Find the order size for which the average cost is exactly \$9 per shirt.
        \textbf{(e)} Can the average cost ever be exactly \$7? Answer using the horizontal asymptote.
        \rubricbox{\gradedblank[id=qd2, type=rubric, points=6, rubric={../rubrics/alg2-u04-practice-test-qd2.yaml}, display={(a) It is a quotient of two polynomials, $7n+180$ over $n$. Algebraically the only
        restriction is $n\ne 0$; in context $n$ must be a positive whole number of shirts, so the
        sensible domain is $n=1,2,3,\dots$ --- an order of $0$ shirts has no ``per shirt'' cost at
        all.\\[3pt]
        (b) $A(20)=\dfrac{140+180}{20}=\dfrac{320}{20}=\$16$ per shirt;
        $A(60)=\dfrac{420+180}{60}=\dfrac{600}{60}=\$10$ per shirt. The average cost is
        \emph{falling} as the order grows, because the one-time \$180 is being split among more
        shirts.\\[3pt]
        (c) Divide term by term: $A(n)=\dfrac{7n}{n}+\dfrac{180}{n}=7+\dfrac{180}{n}$. As
        $n\to\infty$ the second piece shrinks toward $0$, so the horizontal asymptote is $y=7$. That
        \$7 is the true printing cost of one shirt --- the setup fee per shirt is the whole
        difference between the average cost and \$7.\\[3pt]
        (d) Solve $\dfrac{7n+180}{n}=9$. Multiply by $n$ (allowed, since $n\ne0$): $7n+180=9n$, so
        $180=2n$ and $n=90$ shirts.\\[3pt]
        (e) No. $A(n)=7+\dfrac{180}{n}$, and $\dfrac{180}{n}$ is positive for every positive $n$, so
        $A(n)>7$ always. The graph approaches the line $y=7$ but never reaches it --- the setup fee
        never fully disappears, it only gets spread thinner.}]{\linewidth}{5.5cm-6pt}}
\end{enumerate}

\end{document}
